\section{Antennengewinn in dBi und dBd}
\label{section:antennengewinn}
\begin{frame}{Isotropstrahler und Richtwirkung}
    \begin{columns}
        \begin{column}{0.48\linewidth}
            \textbf{Isotropstrahler}: 
            \begin{itemize}
                \item Hypothetische Antenne, die in alle Richtungen gleich stark abstrahlt
                \item Keine reale Antenne, dient als Referenz
            \end{itemize}
            \textbf{Reale Antenne}: 
            \begin{itemize}
                \item Eine reale Antenne weist eine Richtwirkung auf
                \item In bestimmten Richtungen stärker als der Isotropstrahler, in bestimmten Richtungen schwächer als der Isotropstrahler
                \item Die \emph{Hauptstrahlrichtung} ist die Richtung mit dem maximalen Antennengewinn
            \end{itemize}
        \end{column}
        \begin{column}{0.48\linewidth}
            \DARCimage{\linewidth}{751include}
        \end{column}
    \end{columns}
\end{frame}

\begin{frame}
    \frametitle{Gewinn in dBi}
    \begin{itemize}
        \item Gewinn in eine Richtung gegenüber dem Isotropstrahler
        \item Bei Bezug auf den Isotropstrahler wird \textcolor{DARCblue}{dBi} verwendet
    \end{itemize}
\end{frame}

\begin{frame}
    \only<1>{
        \begin{QQuestion}{EG220}{Der Gewinn von Antennen wird häufig in dBi angegeben. Auf welche Vergleichsantenne bezieht man sich dabei? Man bezieht sich dabei auf den~...}{Halbwellenstrahler.}
            {Isotropstrahler.}
            {Horizontalstrahler.}
            {Vertikalstrahler.}
        \end{QQuestion}

    }
    \only<2>{
        \begin{QQuestion}{EG220}{Der Gewinn von Antennen wird häufig in dBi angegeben. Auf welche Vergleichsantenne bezieht man sich dabei? Man bezieht sich dabei auf den~...}{Halbwellenstrahler.}
            {\textbf{\textcolor{DARCgreen}{Isotropstrahler.}}}
            {Horizontalstrahler.}
            {Vertikalstrahler.}
        \end{QQuestion}

    }
\end{frame}

\begin{frame}{Gewinn eines Halbwellendipols}
    \begin{itemize}
        \item Ein Halbwellendipol strahlt senkrecht zum Leiter um \qty{2,15}{\dB} stärker ab als ein Isotropstrahler
        \item Der Gewinn beträgt \qty{2,15}{\dBi}
        \item Bei Bezug auf den Halbwellendipol wird \textcolor{DARCblue}{dBd} verwendet
        \item Ein Halbwellendipol hat in Hauptstrahlrichtung einen Gewinn von \qty{0}{dBd} und \qty{2,15}{\dBi}
    \end{itemize}
\end{frame}

\begin{frame}
    \only<1>{
        \begin{QQuestion}{EG221}{Ein Antennenhersteller gibt den Gewinn einer Antenne mit 5 dBd an. Wie groß ist der Gewinn der Antenne in dBi?}{\qty{2,85}{\dBi}}
            {\qty{5}{\dBi}}
            {\qty{2,5}{\dBi}}
            {\qty{7,15}{\dBi}}
        \end{QQuestion}

    }
    \only<2>{
        \begin{QQuestion}{EG221}{Ein Antennenhersteller gibt den Gewinn einer Antenne mit 5 dBd an. Wie groß ist der Gewinn der Antenne in dBi?}{\qty{2,85}{\dBi}}
            {\qty{5}{\dBi}}
            {\qty{2,5}{\dBi}}
            {\textbf{\textcolor{DARCgreen}{\qty{7,15}{\dBi}}}}
        \end{QQuestion}

    }
\end{frame}%ENDCONTENT
